% !TEX root = ../main.tex

\chapter{Android Client Evaluation Reproducibility}
\label{app:android-evaluation}

\section{Environment and Protocol}

The Android client measurements used an Android 14/API 34 \texttt{sdk\_gphone64\_\allowbreak x86\_64} emulator with four virtual processors and a 192~MiB application heap.
The live-backend trace used the emulator's host-network route to a local backend.
The authentication token was supplied to the instrumentation process at run time and was not retained in an artifact.
The evaluation manifest records the Android build and measurement configuration.
The client results in Section~\ref{sec:android-performance} apply to that configuration.

Startup measurement discarded three cold/foreground launch pairs, then retained 20 cold process starts and 20 task foreground/resume operations.
All retained operations returned a non-empty \texttt{WaitTime} field for the requested activity; the startup runner did not inspect rendered UI content.
Android classified nine as hot starts and returned \texttt{UNKNOWN (0)} for eleven, so their latency distribution uses the retained wait-time field and is not treated as a 20-sample confirmed hot-start distribution.
Controlled state, graph, and event tests used separate warm-up iterations that were not included in the reported samples.
The graph experiment used deterministic payloads with 1, 12, 25, and 50 nodes.
The live event loop performed five repetitions.

\subsection{Factorial Resource Protocol}

The resource-controlled experiment used the same Android~14/API~34 Google
APIs x86\_64 image and Pixel~6 profile in all sessions.
The screen was fixed at \(1080\times2400\) pixels and 420~dpi, and the
application heap remained 228~MiB.
The four configurations combined two or four virtual processor cores with
2{,}048 or 6{,}144~MiB of emulator RAM.
Four counterbalanced blocks changed the order of the configurations.
Each of the sixteen sessions began from a wiped-data cold boot and produced a
manifest, environment record, test output, emulator log, and measurement CSV
files.

The independent unit is one emulator session.
Each session produced 810 repeated controlled samples, and the analysis first
reduced each metric to one median per session.
The four independent values in each factorial cell were then summarized by
their median and quartiles.
Finite bootstrap intervals enumerate all resamples of the four session
summaries.
CPU and RAM effects are block-level percentage changes between the two factor
levels and are interpreted descriptively.
The experiment excludes backend, retrieval, provider, and recommendation
latency.

\section{Metric Definitions}

\begin{longtable}{P{3.7cm}P{10.4cm}}
\caption{Android client metric definitions.}
\label{tab:android-metric-definitions}\\
\toprule
\MnemeTableHeader{Metric} & \MnemeTableHeader{Definition} \\
\midrule
Cold process start &
Elapsed wait time from a forced process start until Android reports completion of the requested activity launch.
The measurement does not inspect rendered UI content. \\
\addlinespace
Task foreground/resume &
Elapsed wait time from requesting the existing task in the foreground.
Success requires a non-empty activity wait-time record; Android's hot-start classifier is recorded separately. \\
\addlinespace
Controlled state transition &
Elapsed time from delivering a deterministic fixture result to the expected ViewModel state.
This excludes live network time and display-frame latency. \\
\addlinespace
Three-stage polling &
Elapsed time until a controlled asynchronous job advances through three client-visible stages and reaches ready state under the production one-second polling schedule. \\
\addlinespace
Graph DOM-ready &
Elapsed time from loading or updating the local graph page until the expected
DOM nodes and renderer marker are observable.
The new-WebView path includes Compose test-rule idle synchronization and can
therefore be affected by an active force animation.
The metric does not establish force-layout convergence. \\
\addlinespace
Selection callback &
Elapsed time from selecting a graph node to receipt of the validated identifier in the native Android layer. \\
\addlinespace
Event stage latency &
Elapsed time for the named queue, retry, replay, upload, or readback operation.
Controlled and live stages are summarized separately. \\
\addlinespace
Median and p95 &
The median is the sample median.
The p95 is the nearest-rank 95th percentile, so the maximum is reported for a five-sample live stage. \\
\bottomrule
\end{longtable}

\section{Sample Inventory}

\begin{longtable}{P{4.2cm}P{1.8cm}P{2.7cm}P{3.7cm}}
\caption{Sample inventory for the Android client evaluation.}
\label{tab:android-sample-inventory}\\
\toprule
\MnemeTableHeader{Group} & \MnemeTableHeader{Samples} & \MnemeTableHeader{Conditions met} & \MnemeTableHeader{Included conditions} \\
\midrule
Startup and resume & 40 & 40 & 20 cold starts; 20 foreground/resume operations \\
State and polling & 160 & 160 & five state conditions at 30 trials; ten polling trials \\
Graph timing & 400 & 400 & ready and selection; four sizes; new and existing WebView \\
Graph-state disclosure & 90 & 90 & ready, fallback, and empty; 30 trials each \\
Controlled event sync & 120 & 120 & queue, failure, retry, and replay; 30 trials each \\
\midrule
Controlled total & 810 & 810 & five result groups reported separately in Chapter~4 \\
\addlinespace
Live-backend stages & 25 & 25 & five stages in each of five repetitions \\
Existing graph/app-flow tests & 11 & 11 & includes Open Paper and restored selection \\
\bottomrule
\end{longtable}

\begin{table}[htbp]
\centering
\caption{Sample inventory for the Android resource-controlled experiment.}
\label{tab:android-factorial-inventory}
\begin{tabular}{P{2.8cm}P{2.0cm}P{2.1cm}P{2.7cm}P{3.0cm}}
\toprule
\MnemeTableHeader{Configuration} & \MnemeTableHeader{Independent sessions} & \MnemeTableHeader{Controlled samples} & \MnemeTableHeader{Successful samples} & \MnemeTableHeader{Functional tests in block 1} \\
\midrule
2 cores / 2~GiB & 4 & 3{,}240 & 3{,}240 & 11 of 11 \\
2 cores / 6~GiB & 4 & 3{,}240 & 3{,}240 & 11 of 11 \\
4 cores / 2~GiB & 4 & 3{,}240 & 3{,}240 & 11 of 11 \\
4 cores / 6~GiB & 4 & 3{,}240 & 3{,}240 & 11 of 11 \\
\midrule
Total & 16 & 12{,}960 & 12{,}960 & 44 of 44 \\
\bottomrule
\end{tabular}
\end{table}
\FloatBarrier

\section{Success Conditions and Retained Artifacts}

A cold or foreground/resume launch succeeds when \texttt{am start -W} returns a non-empty wait-time record for the requested activity.
This condition does not inspect rendered UI content.
A content-state trial succeeds when the ViewModel exposes the expected content origin or retryable error.
A polling trial succeeds when all three processing stages are observed in order and the paper reaches ready state.
A graph trial succeeds when the page reports ready and the selected identifier reaches the native client.
Ready, fallback, and empty tests additionally require the corresponding reader-visible disclosure.

An event trial succeeds when the local queue and reservation state match the expected stage.
After a retryable failure, the batch must return to pending state.
After success, accepted and duplicate counts must account for the full batch.
For the live trace, each repetition must queue three events, recover from the injected failure, upload three new identifiers, recognize all three on replay, and return a live-origin briefing to the client.
Briefing entry count and preference-model version are recorded as outcomes and do not determine trial success.

The retained evaluation package contains direct comma-separated measurements, an environment record, a run manifest without the authentication token, connected Android test XML, machine-readable summaries, and the vector figures used in Chapter~4.
The five live briefing responses contain zero entries and retain preference-model version 1.
Those values remain part of the reproducibility record because they bound the event-loop conclusion.
